
\subsection{The Galois group}
The Galois group is a group of automorphisms of a field extension 
that fixes the base field. The Galois group is a central concept in 
Galois theory, since it's precisley what the theory is built upon.

\begin{definition}[Automorphism]
A field-automorphism for the field $(F, +, \cdot)$ is a bijective map 
$\phi: F \to F$ such that for all $a, b \in F$, the basic properties of a field are preserved, 
meaning that $\phi$ act as a linear mapping with respect to the ring operations:

$$ \phi(a + b) = \phi(a) + \phi(b) $$
$$ \phi(ab) = \phi(a)\phi(b) $$ \vspace{0.2em}
\end{definition}

\begin{lemma}[Zero-lemma]
For any automorphism $\phi$ of a field $(F, +, \cdot)$, we have that:
\end{lemma}
\begin{proof}
$$ \phi(0) = \phi(0 \cdot a) = \phi(0) \phi(a) $$
$$ \therefore \phi(0) = 0 $$
\end{proof}


\begin{lemma}[Unit-lemma]
For any automorphism $\phi$ of a field $(F, +, \cdot)$, we have that:
\end{lemma}
\begin{proof}
$$ \phi(a) = \phi(1 \cdot a) = \phi(1) \phi(a) $$
$$ \therefore \phi(1) = 1 $$
\end{proof}

\begin{lemma}
    All the automorphisms of a given field form a group under the operation of composition.
\end{lemma}
\begin{proof}
    Consider two automorphisms $\phi_1$ and $\phi_2$ of a field $F$. The composition
    $\phi_1 \circ \phi_2$ is also an automorphism of $F$, since it preserves the field operations:

    $$ (\phi_1 \circ \phi_2)(a + b) = \phi_1(\phi_2(a + b)) = \phi_1(\phi_2(a) + \phi_2(b)) = \phi_1(\phi_2(a)) + \phi_1(\phi_2(b)) $$
    $$ (\phi_1 \circ \phi_2)(ab) = \phi_1(\phi_2(ab)) = \phi_1(\phi_2(a) \cdot \phi_2(b)) = \phi_1(\phi_2(a)) \cdot \phi_1(\phi_2(b)) $$
    $$ \therefore \phi_1 \circ \phi_2 \texttt{ is an automorphism of } F $$
\end{proof}

\begin{definition}[Galois group]
    The Galois group of a field extension $E$ for the field $F$ is the group of
    automorphisms of $E$ that fix $F$. In other words, the Galois group is the
    group of all the automorphisms of $E$ that leave every element of $F$ unchanged.

    $$ Aut(E) = \{ \phi : E \rightarrow E \texttt{ automorphisms } \} $$
    $$ Gal_{F}(E) = \{ \phi \in \text{Aut}(E) \;:\; \forall a \in F \;\; \phi(a) = a \} $$
\end{definition} \vspace{0.2em}

\begin{lemma}[roots-permutation lemma]
    If $\phi$ is an automorphism of a field extension $E$ of $F$, meaning it's an element 
    of the Galois group $Gal_{F}(E)$, then, for every polynomal $p \in F[x]$ the automorphism 
    $\phi$ must map roots of $p$ to other roots of $p$. 
\end{lemma}
\begin{proof}
    $$ \phi(p(x)) = \phi(a_nx^n + a_{n-1}x^{n-1} + \dots + a_0) $$
    $$ = a_n \phi(x^n) + a_{n-1} \phi(x^{n-1}) + \dots + a_0 $$
    $$ = a_n \phi(x)^n + a_{n-1} \phi(x)^{n-1} + \dots + a_0 $$
    $$ = p(\phi(x)) $$
    $$ \therefore \phi(p(x)) = p(\phi(x)) $$
\end{proof}   

The roots-permutation lemma is important, since it highlights a very intresting property of the
Galois group, which is the fact that for every polynomial $p \in F[x]$ that has roots 
$r_1, r_2, \dots, r_n$ not contained in $F$, the Galois group $Gal_{F}(E)$ will permute 
the roots of $p$ and will herby be isomorphic to the group of permutations of $n$ elements $S_n$.

\begin{theorem}[Fundamental theorem of Galois theory]
    Take a field extension $E$ of a field $F$, then there exists a bijection between 
    the intermediate fields $F \subseteq I \subseteq E$ and the subgroups of the
    Galois group $Gal_{F}(E)$, defined as follows:

    $$ \mathcal{G} : \{ I \;:\; F \subseteq I \subseteq E \} \rightarrow \texttt{ subgroups of } Gal_{F}(E) $$
    $$ \mathcal{G}(I) = Gal_{I}(E)  $$
\end{theorem}
\begin{proof}
    The function $\mathcal{G}$ is clearly injective, so we only need to show that it's surjective. To 
    show that, we will consider the generic subgroup $H$ of the Galois group 
    $Gal_{F}(E)$, and the set $I$ of all the elements of $E$ that are fixed by every
    automorphism in $H$. We will show that $I$ is in the domain of $\mathcal{G}$ such that
    $\mathcal{G}(I) = H$ by its very construction. \\

    First, let's show that not only every element of $I$ is fixed by every
    automorphism in $H$, but also that every product and sum of elements of $I$ is 
    still fixed by every automorphism in $H$, herby showing that $I$ is closed.

    $$ \forall a,b \in I \;\; \forall \phi \in H \;\; \phi(a \cdot b) = \phi(a)\phi(b) = ab $$
    $$ \forall a,b \in I \;\; \forall \phi \in H \;\; \phi(a + b) = \phi(a) + \phi(b) = a + b $$ \vspace{0.2em}

    Since $H$ is a subgroup of $Gal_{F}(E)$, we have that every $\phi$ in $H$ fixes $F$, meaning that
    $F \subseteq I$. Thus $I$ is a field extension of $F$ that maps to $H$ via $\mathcal{G}$ (Galois correspondence). \\
\end{proof}

A very intresting consequence of the fundamental theorem of Galois theory is that 
every subgroup of a Galois group $Gal_{F}(E)$ is still a Galois group of some smaller field
extension.

\newpage

\begin{theorem}
\label{thm:GaloisIsomorphism}
    As a consequence of the first theorem of isomorphism (theorem \ref{thm:firstIsomorphism}) 
    we can claim that there is an isomorphism between $Gal_{A}(B)$ and $Gal_{A}(S)/Gal_{B}(S)$ 
    given that $A \subseteq B \subseteq S$ and that $B$ is actually a \emph{normal extension} 
    of the field $A$.
\end{theorem}
\begin{proof}
    Let's define an homomorphism $\phi$ from $Gal_{B}(S)$ to $Gal_{B}(A)$ as shown below:

    $$ \psi : Gal_{A}(S) \rightarrow Gal_{A}(B) $$
    $$ \psi(\sigma : S \rightarrow S) = \sigma : B \rightarrow B $$ 
    \vspace{0.1em}

    This construction of $\psi$ makes it map an automorphism that fixes $A$ in $S$ to 
    an automorphism that is still fixing $A$ but restricted to $B$. We need to show that it's 
    an epimorphism, meaning that it is an homomorphism and that it's surjective, 
    to then be able to inspect it's kernel and apply the first theorem of isomorphism 
    (theorem \ref{thm:firstIsomorphism}). \\

    To show that $\psi$ is an epimorphism, we start by showing that it is a homomorphism:

    \vspace{0.2em}
    $$ \psi(\sigma \circ \pi) = \sigma \circ \pi = \psi(\sigma) \circ \psi(\pi) $$
    \vspace{0.2em}

    Since $\psi$ is compatible with the composition operation, we can say that it is a well-defined 
    homomorphism. Now we need to show that it is surjective. The normality of $B$ as an extension
    of $A$ is crucial to establish surjectivity here, because $B$ is normal, $B$ is the splitting 
    field of some polynomial over $A$. By the \textbf{extension theorem for automorphisms}, every 
    automorphism $\tau \in \text{Gal}_A(B)$ extends to some $\tilde{\tau} \in 
    \text{Gal}_A(S)$. \\  

    $$
        \forall \tau \in \text{Gal}_A(B), \; \exists \tilde{\tau} \in \text{Gal}_A(S) 
        \text{ such that } \tilde{\tau}|_B = \tau.
    $$
    \vspace{0.2em}
    
    Thus, $\psi(\tilde{\tau}) = \tau$, proving $\psi$ is surjective. If $B$ wasn't normal in $A$, 
    there would exist automorphisms of $B$ over $A$ that cannot extend to $S$, 
    breaking surjectivity. For example, if $B$ lacked a root required to 
    define $\tilde{\tau}$, 
    $\psi$ would fail to map onto $\text{Gal}_A(B)$. \\
    
    Established that we're working with an epimorphism, we can now inspect its kernel:

    $$ \; [ \; \psi(\sigma \circ \gamma) = \sigma \; ] \; \iff \forall b \in B\;\; \gamma(b) = b $$
    $$ \texttt{ker}_{\psi} = Gal_{B}(S) $$
    \vspace{0.2em}

    We can apply theorem \ref{thm:firstIsomorphism} (first theorem of isomorphism) 
    to the current setup obtaining:

    $$ Gal_{A}(B) = \vv{\psi}(Gal_{B}(S)) 
    \sim \frac{Gal_{A}(S)}{\texttt{ker}_{\psi}} 
    = \frac{Gal_{A}(S)}{Gal_{B}(S)} $$

\end{proof}

\newpage